\documentclass[a4paper]{article}
\usepackage{amsmath, amsthm, amssymb, bm, graphicx, hyperref, mathrsfs}
\usepackage{booktabs}
\usepackage[affil-it]{authblk}
\usepackage[backend=bibtex,style=numeric]{biblatex}

\usepackage{geometry}
\geometry{margin=1.5cm, vmargin={0pt,1cm}}
\setlength{\topmargin}{-1cm}
\setlength{\paperheight}{29.7cm}
\setlength{\textheight}{25.3cm}

\addbibresource{citation.bib}

\begin{document}
% =================================================
\title{Numerical Analysis Homework 2}

\author{Yilin Chen 3230101092
  \thanks{Electronic address: \texttt{yilinchen@zju.edu.cn}}}
\affil{(Numerical Analysis, Math2301), Zhejiang University }


\date{Due time:2025/10/27 \today}

\maketitle

% \begin{abstract}
%     The abstract is not necessary for the theoretical homework, 
%     but for the programming project, 
%     you are encouraged to write one.      
% \end{abstract}





% ============================================
\section*{I. Linear Interpolation Error}

\subsection*{I-a}

According to Lagrange formula, $p_1(f;x) = f(x_0)\dfrac{x-x_1}{x_0 - x_1} + f(x_1)\dfrac{x-x_0}{x_1 - x_0}$, so 

$$f''(\xi(x)) = \dfrac2{\xi^3(x)} = 2\dfrac{f(x) - p_1(f;x)}{(x - x_0)(x - x_1)} = 2\dfrac{\frac1x + \frac{x}2- \frac32}{(x - 1)(x - 2)} = \dfrac{1}{x}$$

Thus $\xi(x) = \sqrt[3]{2x}$


\subsection*{I-b}

Because $\xi(x) = \sqrt[3]{2x}$ and $f''(\xi(x)) = \dfrac1x$, we know the 2 function are increasing and decreasing on$[1, 2]$ respectively, thus

$$\max\xi(x) = \xi(2) = \sqrt[3]{4}$$

$$\min\xi(x) = \xi(1) = \sqrt[3]{2}$$

$$\max f''(\xi(x)) = \dfrac11 = 1$$

\section*{II. Non-negative polynomial interpolation}

Imitate the construction of Lagrange formula, let $$p = \sum_{i = 1}^{n}f_i \left(\dfrac{\prod\limits_{j \neq i, j = 0}^{n}(x - x_j)}{\prod\limits_{j \neq i, j = 0}^{n}(x_i - x_j)}\right)^2$$

Obviously, p satisfies the requirement.

\section*{III. Divided differences of exponential function}

\subsection*{III-a}

\textbf{Base case:} For \( n = 0 \), 
\[
f[t] = e^t = \frac{(e-1)^0}{0!} e^t.
\]
For \( n = 1 \),
\[
f[t, t+1] = \frac{e^{t+1} - e^t}{(t+1) - t} = e^t(e - 1) = \frac{(e-1)^1}{1!} e^t.
\]

\textbf{Inductive step:} Assume the formula holds for \( n \),then 
\[
f[t+1, \ldots, t+n+1] = \frac{(e-1)^n}{n!} e^{t+1}, \quad f[t, \ldots, t+n] = \frac{(e-1)^n}{n!} e^t.
\]
From theorem 2.17:
\[
f[t, \ldots, t+n+1] = \frac{f[t+1, \ldots, t+n+1] - f[t, \ldots, t+n]}{(t+n+1) - t}.
\]
Thus,
\[
f[t, \ldots, t+n+1] = \frac{(e-1)^n}{n!} \cdot \frac{e^{t+1} - e^t}{n+1} = \frac{(e-1)^n}{n!} \cdot \frac{e^t(e - 1)}{n+1} = \frac{(e-1)^{n+1}}{(n+1)!} e^t.
\]
This completes the induction.

\subsection*{III-b}

From Corollary 2.22, there exists \( \xi \in (0, n) \) such that
\[
f[0, 1, \ldots, n] = \frac{1}{n!} f^{(n)}(\xi).
\]
Since \( f^{(n)}(x) = e^x \), we have
\[
f[0, 1, \ldots, n] = \frac{1}{n!} e^\xi.
\]
From the induction result with \( t = 0 \),
\[
f[0, 1, \ldots, n] = \frac{(e - 1)^n}{n!}.
\]
Equating the two expressions:
\[
\frac{(e - 1)^n}{n!} = \frac{1}{n!} e^\xi \quad \Rightarrow \quad e^\xi = (e - 1)^n \quad \Rightarrow \quad \xi = n \ln(e - 1).
\]

Now, compare \( \xi = n \ln(e - 1) \) with the midpoint \( n/2 \). Since \( e - 1 \approx 1.71828 \) and \( \ln(e - 1) \approx 0.54132 \), we have
\[
\xi \approx 0.54132n > 0.5n = \frac{n}{2}.
\]
Therefore, \( \xi \) is located to the right of the midpoint \( n/2 \).

\section*{IV. Interpolation and Minimum Location}

\subsection*{IV-a}

Given the data points: \( f(0) = 5 \), \( f(1) = 3 \), \( f(3) = 5 \), \( f(4) = 12 \), we construct the divided difference table:

\begin{center}
\begin{tabular}{c|cccc}
$x$ & $f[x]$ &  & &  \\
\hline
0 & 5 \\
1 & 3 & -2 \\
3 & 5 & 1 & 1 \\
4 & 12 & 7 & 2 & 0.25 \\
\end{tabular}
\end{center}

The Newton interpolation polynomial is:
\[
p_3(x) = 5 + (-2)(x-0) + 1 \cdot (x-0)(x-1) + 0.25 \cdot (x-0)(x-1)(x-3) = 0.25x^3 - 2.25x + 5
\]

\subsection*{IV-b}

The data suggest a minimum in \( (1, 3) \). We find the critical point of \( p_3(x) \):
\[
p_3'(x) = 0.75x^2 - 2.25.
\]
Setting \( p_3'(x) = 0 \):
\[
0.75x^2 - 2.25 = 0 \quad \Rightarrow \quad x^2 = 3 \quad \Rightarrow \quad x = \sqrt{3}.
\]
The second derivative test confirms this is a minimum:
\[
p_3''(x) = 1.5x > 0 \quad \text{for } x > 0.
\]

Therefore, the approximate location of the minimum is \( x_{\min} = \sqrt{3}\).

\section*{V. Divided Difference of \( f(x) = x^7 \)}

\subsection*{V-a}

We compute the divided difference using the recursive definition and the fact that for repeated nodes, the divided difference involves derivatives.

Given \( f(x) = x^7 \), we have:
\[
f(0) = 0, \quad f(1) = 1, \quad f(2) = 128.
\]
The derivatives are:
\[
f'(x) = 7x^6, \quad f''(x) = 42x^5.
\]

From Corollary 2.23, $f[1, 1] = \dfrac{f'(1)}{1!} = 7$, $f[2, 2] = \dfrac{f'(2)}{1!} = 448$, $f[1, 1, 1] = \dfrac{f''(1)}{2!} = 21$

We construct the divided difference table for nodes \( 0, 1, 1, 1, 2, 2 \):

\begin{center}
\begin{tabular}{c|cccccc}
$x$ & $f[\cdot]$ & & & & & \\
\hline
0 & 0 \\
1 & 1 & 1 \\
1 & 1 & 7 & 6 \\
1 & 1 & 7 & 21 & 15\\
2 & 128 & 127 & 120 & 99 & 42\\
2 & 128 & 448 & 321 & 201 & 102 \\
\end{tabular}
\end{center}

That is, \( f[0,1,1,1,2,2] = 30 \).

\subsection*{V-b}

By Corollary 2.22, there exists \( \xi \in (0,2) \) such that
\[
f[0,1,1,1,2,2] = \frac{f^{(5)}(\xi)}{5!}.
\]
We compute \( f^{(5)}(x) = 7 \cdot 6 \cdot 5 \cdot 4 \cdot 3 \cdot x^2 = 2520 x^2 \). Thus,
\[
\frac{2520 \xi^2}{120} = 21 \xi^2 = 30 \quad \Rightarrow \quad \xi^2 = \frac{30}{21} = \frac{10}{7}.
\]
Hence, \( \xi = \sqrt{\frac{10}{7}} \) (since \( \xi > 0 \)).

\section*{VI. Hermite Interpolation and Error Estimation}

\subsection*{VI-a}

Given the conditions:
\[
f(0) = 1, \quad f(1) = 2, \quad f'(1) = -1, \quad f(3) = 0, \quad f'(3) = 0,
\]
we construct the Hermite interpolation polynomial \( H(x) \) of degree at most 4 using the nodes \( 0, 1, 1, 3, 3 \). The divided difference table is:

\begin{center}
\begin{tabular}{c|ccccc}
$x$ & $f[\cdot]$ & & & & \\
\hline
0 & 1 \\
1 & 2 & 1 \\
1 & 2 & -1 & -2 \\
3 & 0 & -1 & 0 & $\frac23$ \\
3 & 0 & 0 & $\frac12$ & $\frac14$ & $-\frac{5}{36}$\\
\end{tabular}
\end{center}

The polynomial is:
\[
H(x) = 1 + 1 \cdot x + (-2) \cdot x(x-1) + \frac{2}{3} \cdot x(x-1)^2 + \left(-\frac{5}{36}\right) \cdot x(x-1)^2(x-3).
\]

Evaluating at \( x = 2 \):

$$H(2) = 1 + 2 + (-2)(2)(1) + \frac{2}{3}(2)(1)^2 + \left(-\frac{5}{36}\right)(2)(1)^2(-1) = \frac{11}{18}$$


\subsection*{VI-b}

Given \( f \in C^5[0, 3] \) and \( |f^{(5)}(x)| \leq M \) on \([0, 3]\), the error formula for Hermite interpolation is:
\[
f(x) - H(x) = \frac{f^{(5)}(\xi)}{5!} \cdot (x - 0)(x - 1)^2(x - 3)^2, \quad \xi \in (0, 3).
\]

At \( x = 2 \):
\[
(2 - 0)(2 - 1)^2(2 - 3)^2 = 2 \cdot 1 \cdot 1 = 2,
\]
so the error is:
\[
|f(2) - H(2)| = \left| \frac{f^{(5)}(\xi)}{120} \cdot 2 \right| = \frac{|f^{(5)}(\xi)|}{60} \leq \frac{M}{60}.
\]

Thus, the maximum possible error is:
\[
\frac{M}{60}.
\]

\section*{VII. Forward and Backward Differences}

\subsection*{VII-a}

We prove by induction that
\[
\Delta^k f(x) = k! h^k f[x_0, x_1, \ldots, x_k],
\]
where \( x_j = x + jh \).

\textbf{Base case:} For \( k = 1 \),
\[
\Delta f(x) = f(x+h) - f(x) = \frac{f(x_1) - f(x_0)}{x_1 - x_0} \cdot (x_1 - x_0) = f[x_0, x_1] \cdot h = 1! h^1 f[x_0, x_1].
\]

\textbf{Inductive step:} Assume the identity holds for some \( k \). Then
\[
\Delta^{k+1} f(x) = \Delta^k f(x+h) - \Delta^k f(x).
\]
By the inductive hypothesis,
\[
\Delta^k f(x+h) = k! h^k f[x_1, x_2, \ldots, x_{k+1}], \quad \Delta^k f(x) = k! h^k f[x_0, x_1, \ldots, x_k].
\]
Thus,
\[
\Delta^{k+1} f(x) = k! h^k \left( f[x_1, \ldots, x_{k+1}] - f[x_0, \ldots, x_k] \right).
\]
From the property of divided differences,
\[
f[x_0, x_1, \ldots, x_{k+1}] = \frac{f[x_1, \ldots, x_{k+1}] - f[x_0, \ldots, x_k]}{x_{k+1} - x_0},
\]
and since \( x_{k+1} - x_0 = (k+1)h \), we have
\[
f[x_1, \ldots, x_{k+1}] - f[x_0, \ldots, x_k] = (k+1)h f[x_0, x_1, \ldots, x_{k+1}].
\]
Therefore,
\[
\Delta^{k+1} f(x) = k! h^k \cdot (k+1)h f[x_0, x_1, \ldots, x_{k+1}] = (k+1)! h^{k+1} f[x_0, x_1, \ldots, x_{k+1}].
\]
This completes the induction.

\subsection*{VII-b}

We prove by induction that
\[
\nabla^k f(x) = k! h^k f[x_0, x_{-1}, \ldots, x_{-k}],
\]
where \( x_j = x + jh \) (with negative indices).

\textbf{Base case:} For \( k = 1 \),
\[
\nabla f(x) = f(x) - f(x-h) = \frac{f(x_0) - f(x_{-1})}{x_0 - x_{-1}} \cdot (x_0 - x_{-1}) = f[x_{-1}, x_0] \cdot h.
\]
By symmetry of divided differences, \( f[x_{-1}, x_0] = f[x_0, x_{-1}] \), so
\[
\nabla f(x) = h f[x_0, x_{-1}] = 1! h^1 f[x_0, x_{-1}].
\]

\textbf{Inductive step:} Assume the identity holds for some \( k \). Then
\[
\nabla^{k+1} f(x) = \nabla^k f(x) - \nabla^k f(x-h).
\]
By the inductive hypothesis,
\[
\nabla^k f(x) = k! h^k f[x_0, x_{-1}, \ldots, x_{-k}], \quad \nabla^k f(x-h) = k! h^k f[x_{-1}, x_{-2}, \ldots, x_{-(k+1)}].
\]
Thus,
\[
\nabla^{k+1} f(x) = k! h^k \left( f[x_0, x_{-1}, \ldots, x_{-k}] - f[x_{-1}, x_{-2}, \ldots, x_{-(k+1)}] \right).
\]
From the property of divided differences,
\[
f[x_0, x_{-1}, \ldots, x_{-(k+1)}] = \frac{f[x_0, x_{-1}, \ldots, x_{-k}] - f[x_{-1}, x_{-2}, \ldots, x_{-(k+1)}]}{x_0 - x_{-(k+1)}},
\]
and since \( x_0 - x_{-(k+1)} = (k+1)h \), we have
\[
f[x_0, x_{-1}, \ldots, x_{-k}] - f[x_{-1}, x_{-2}, \ldots, x_{-(k+1)}] = (k+1)h f[x_0, x_{-1}, \ldots, x_{-(k+1)}].
\]
Therefore,
\[
\nabla^{k+1} f(x) = k! h^k \cdot (k+1)h f[x_0, x_{-1}, \ldots, x_{-(k+1)}] = (k+1)! h^{k+1} f[x_0, x_{-1}, \ldots, x_{-(k+1)}].
\]
This completes the induction.

\section*{VIII. $\dfrac\partial{\partial x_0}f[x_0, x_1, \cdots, x_n]$}

We prove by induction that

$$\dfrac\partial{\partial x_0}f[x_0, x_1, \cdots, x_n] = f[x_0, x_0, x_1, \cdots, x_n]$$

\textbf{Base case:}

$$\dfrac\partial{\partial x_0}f[x_0] = \dfrac{\mathrm{d}} {\mathrm{d}  x_0}f(x_0) = f'(x_0) = f[x_0, x_0]$$

\textbf{Inductive step:} Assume the identity holds for $n = k$. Then for $n = k+1$,

\begin{align*}
\dfrac\partial{\partial x_0}f[x_0, x_1, \cdots, x_{k+1}] &= \dfrac\partial{\partial x_0}\dfrac{f[x_0, \cdots, x_{k-1}, x_{k+1}] - f[x_0, \cdots, x_{k-1}, x_{k}]}{x_{k+1} - x_k}\\
&= \dfrac{f[x_0, x_0, x_1, \cdots, x_{k-1}, x_{k+1}] - f[x_0, x_0, x_1, \cdots, x_{k-1}, x_{k}]}{x_{k+1} - x_k}\\
&= f[x_0, x_0, x_1, \cdots, x_k, x_{k+1}]
\end{align*}

This completes the induction.  

\section*{IX. Min-Max Problem}

We wish to minimize
\[
\max_{x \in [a, b]} \left| a_0 x^n + a_1 x^{n-1} + \cdots + a_n \right|
\]
over all choices of \( a_1, a_2, \ldots, a_n \in \mathbb{R} \), with \( a_0 \neq 0 \) fixed.

This is equivalent to finding the monic polynomial of degree \( n \) (after scaling) that minimizes the maximum absolute value on the interval \([a, b]\). The solution is given by the Chebyshev polynomial of degree \( n \) scaled appropriately.

Let \( P(x) = a_0 x^n + a_1 x^{n-1} + \cdots + a_n \). Define the linear transformation mapping \([a, b]\) to \([-1, 1]\):
\[
x = \frac{b - a}{2} t + \frac{a + b}{2}, \quad t \in [-1, 1].
\]
Then
\[
P(x) = a_0 \left( \frac{b - a}{2} t + \frac{a + b}{2} \right)^n + \text{lower degree terms}.
\]
The leading term becomes \( a_0 \left( \frac{b - a}{2} \right)^n t^n \). Define
\[
Q(t) = \frac{P(x)}{a_0} = \left( \frac{b - a}{2} \right)^n t^n + \text{lower degree terms}.
\]
Let \( R(t) = Q(t) / \left( \frac{b - a}{2} \right)^n \). Then \( R(t) \) is a monic polynomial of degree \( n \). From Corollary 2.48, it is known that the monic polynomial of degree \( n \) that minimizes the maximum absolute value on \([-1, 1]\) is \( \frac{1}{2^{n-1}} T_n(t) \), where \( T_n(t) \) is the Chebyshev polynomial of degree \( n \), and the minimum value is \( \frac{1}{2^{n-1}} \). Thus,
\[
\min \max_{t \in [-1, 1]} |R(t)| = \frac{1}{2^{n-1}}.
\]
Then
\[
\min \max_{t \in [-1, 1]} |Q(t)| = \left( \frac{b - a}{2} \right)^n \cdot \frac{1}{2^{n-1}} = \frac{(b - a)^n}{2^{2n - 1}}.
\]
Since \( \max_{x \in [a, b]} |P(x)| = |a_0| \max_{t \in [-1, 1]} |Q(t)| \), we obtain
\[
\min \max_{x \in [a, b]} |P(x)| = \frac{|a_0|(b - a)^n}{2^{2n - 1}}.
\]

\section*{X. Minimax Property of Scaled Chebyshev Polynomial}

Let $x_k' = \cos{\frac{k}{n}\pi}, k = 0, 1, \cdots, n$, then $T_n$ have extreme values at the $(n+1)$ points:

$$T_n(x'_k) = (-1)^k$$

Suppose $||\hat{p_n}||_\infty \le ||p||_\infty$ does not hold. Then Theorem 2.45 implies that

\begin{equation}\label{eq1}
  \exists p \in \mathbb{P}_n^a, s.t. |p(x)| < \dfrac{1}{T_n(a)}
\end{equation}

Let $Q(x) = \dfrac{T_n(x)}{T_n(a)} - p(x)$, then

$$Q(x_k') = \dfrac{(-1)^k}{T_n(a)} - p(x_k')$$

By (2.46), $Q(x)$ has alternating signs at these $n+1$ points. Hence $Q(x)$ must have $n$ zeros. However, by the construction of $Q(x)$, the degree of $Q(x)$ is at most $n - 1$. 

Therefore, $Q(x) \equiv 0$ and $p(x) = \dfrac{T_n(x)}{T_n(a)}$, which implies that

$$\max|p(x)| = \dfrac1{T_n(a)}$$

This is a contradiction to \eqref{eq1}.

\section*{XI. Proof of Lemma 2.53}

According to the definition : $\forall k=0,1, \ldots, n, \quad b_{n, k}(t):=\binom{n}{k} t^k(1-t)^{n-k}$

\begin{align*}
  \dfrac{n-k}{n}b_{n,k}(t) + \dfrac{k+1}{n}b_{n, k+1}(t) &= \dfrac{n-k}{n}\dfrac{n!}{k!(n-k)!}t^k(1-t)^{n-k} + \dfrac{k+1}n \dfrac{n!}{(k+1)!(n-k-1)!}t^{k+1}(1-t)^{n-k-1}\\
  &= \dfrac{(n-1)!}{k!(n-k-1)!}t^k(1-t)^{n-k} + \dfrac{(n-1)!}{k!(n-k-1)!}t^{k+1}(1-t)^{n-k-1}\\
  &= \dfrac{(n-1)!}{k!(n-k-1)!}t^k(1-t)^{n-k-1}(1-t+t)\\
  &= b_{n-1,k}(t)
\end{align*}

\section*{XII. Proof of Lemma 2.55}

\begin{align*}
  \int_0^1 b_{n, k}(t) \mathrm{d}t &= \int_0^1 \binom{n}{k} t^k(1-t)^{n-k} \mathrm{d}t\\
  &= \frac{n!}{k!(n-k)!} B(k+1, n-k+1)\\
  &= \frac{n!}{k!(n-k)!} \frac{\Gamma(k+1)\Gamma(n-k+1)}{\Gamma(n+2)}\\ 
  &= \frac{n!}{k!(n-k)!} \frac{k!(n-k)!}{(n+1)!}\\
  &= \frac1{n+1}
\end{align*}

% ===============================================
% \section*{ \center{\normalsize {Acknowledgement}} }
% Give your acknowledgements here(if any).


% \printbibliography

% If you are not familiar with \texttt{bibtex}, 
% it is acceptable to put a table here for your references.
\end{document}